\documentclass[xcolor=dvipsnames,t,aspectratio=169]{beamer}

\usecolortheme{rose}
\usecolortheme{dolphin}
\usetheme{Boadilla}

\input{../../templates/slides/imports}
\input{../../templates/slides/settings}
\input{../../templates/slides/commands}

\titlegraphic{
    \includegraphics[scale = 0.5]{../../templates/slides/logo}
}

\logo{
\begin{tikzpicture}[overlay,remember picture]
\node[xshift=-.75cm, yshift=-1cm] at (current page.30){
    \includegraphics[width=0.1\textwidth]{../../templates/slides/logo}};
\end{tikzpicture}
}
\usetikzlibrary{positioning}

\newcommand{\highlight}[1]{{\color{nes_dark_orange} #1}}

\title{Feature Engineering, Web Scraping e APIs}

\author{Eduardo Adame}

\date{{\color{nes_dark_purple}\textbf{Ciência de Dados}\\[0.5em] 15 de maio de 2026 }}

\begin{document}

\frame[plain]{\titlepage}
\setcounter{framenumber}{0}


% ============================================================
\section{Feature Engineering}
% ============================================================

\begin{frame}[c]{O que é Feature Engineering?}

\begin{columns}
\column{0.5\textwidth}
\textbf{\color{nes_dark_purple}Definição}
\begin{itemize}
    \item Processo de \highlight{criar, transformar e selecionar} variáveis a partir dos dados brutos
    \item Objetivo: fornecer ao modelo informação mais relevante e na forma certa
    \item Não é sobre o modelo — é sobre \highlight{o que você alimenta} nele
\end{itemize}

\textbf{\color{nes_dark_purple}Por que importa?}
\begin{itemize}
    \item Modelos simples com boas features batem modelos complexos com features ruins
    \item Reduz overfitting ao eliminar ruído
    \item Diminui custo computacional
\end{itemize}

\column{0.5\textwidth}
\textbf{\color{nes_dark_purple}Tipos de transformação}
\begin{itemize}
    \item \highlight{Numéricas}: normalização, log, binning, interações
    \item \highlight{Categóricas}: encoding, target encoding, frequência
    \item \highlight{Temporais}: hora, dia da semana, sazonalidade
    \item \highlight{Texto}: TF-IDF, embeddings, contagem de tokens
    \item \highlight{Derivadas}: razões, diferenças, agregações por grupo
\end{itemize}
\end{columns}

\end{frame}


\begin{frame}[c, fragile]{Transformações numéricas}

\begin{code}[Normalização e transformações]{python}
import pandas as pd
import numpy as np
from sklearn.preprocessing import StandardScaler, MinMaxScaler

# Padronização (média 0, desvio 1) — para modelos lineares e redes neurais
scaler = StandardScaler()
df["preco_std"] = scaler.fit_transform(df[["preco"]])
# Min-Max (intervalo [0, 1]) — para KNN, SVM
scaler_mm = MinMaxScaler()
df["preco_norm"] = scaler_mm.fit_transform(df[["preco"]])
# Log — para distribuições assimétricas (renda, preço de imóveis)
df["preco_log"] = np.log1p(df["preco"])  # log1p = log(1 + x), evita log(0)
# Binning — transformar contínua em ordinal
df["faixa_preco"] = pd.cut(df["preco"],
                           bins=[0, 100, 500, np.inf],
                           labels=["barato", "médio", "caro"])
\end{code}

\end{frame}


\begin{frame}[c, fragile]{Encoding de variáveis categóricas}

\begin{code}[One-Hot, Ordinal e Target Encoding]{python}
# One-Hot Encoding — para categorias sem ordem (cor, cidade)
df = pd.get_dummies(df, columns=["cidade"], drop_first=True)
# Ordinal Encoding — para categorias com ordem (faixa etária, grau)
from sklearn.preprocessing import OrdinalEncoder
oe = OrdinalEncoder(categories=[["baixo", "médio", "alto"]])
df["grau_enc"] = oe.fit_transform(df[["grau"]])
# Target Encoding — substitui categoria pela média do alvo
# (cuidado: vaza informação do alvo se não usar cross-fitting)
media_por_cidade = df.groupby("cidade")["preco"].mean()
df["cidade_target"] = df["cidade"].map(media_por_cidade)
# Frequência — para alta cardinalidade (CEP, ID de produto)
freq = df["produto_id"].value_counts(normalize=True)
df["produto_freq"] = df["produto_id"].map(freq)
\end{code}

\end{frame}


\begin{frame}[c, fragile]{Features temporais e derivadas}

\begin{code}[Extraindo informação de datas e criando interações]{python}
df["data"] = pd.to_datetime(df["data"])

# Decomposição temporal
df["hora"]         = df["data"].dt.hour
df["dia_semana"]   = df["data"].dt.dayofweek   # 0 = segunda
df["mes"]          = df["data"].dt.month
df["fim_de_semana"] = df["dia_semana"].isin([5, 6]).astype(int)

# Sazonalidade cíclica — hora não é linear (23h e 0h são próximas!)
df["hora_sin"] = np.sin(2 * np.pi * df["hora"] / 24)
df["hora_cos"] = np.cos(2 * np.pi * df["hora"] / 24)

# Features derivadas e interações
df["preco_por_m2"]   = df["preco"] / df["area"]
df["idade_km"]       = df["idade_carro"] * df["km_rodados"]  # interação
df["media_por_grupo"] = df.groupby("cidade")["preco"] \
                          .transform("mean")  # agregação sem vazar teste
\end{code}

\end{frame}


% ============================================================
\section{Model Selection e Cross-Validation}
% ============================================================

\begin{frame}[c]{O problema de avaliar modelos}

\begin{columns}
\column{0.5\textwidth}
\textbf{\color{nes_dark_purple}Por que não usar o treino para avaliar?}
\begin{itemize}
    \item O modelo \highlight{memoriza} os dados de treino
    \item Erro de treino $\ll$ erro real — \emph{overfitting}
    \item Precisamos estimar o desempenho em dados \highlight{nunca vistos}
\end{itemize}

\vspace{0.5em}
\textbf{\color{nes_dark_purple}Holdout simples}
\begin{itemize}
    \item Separar treino / teste uma vez
    \item Rápido, mas \highlight{instável}: depende da divisão
    \item Com dados limitados, descarta informação valiosa
\end{itemize}

\column{0.5\textwidth}
\textbf{\color{nes_dark_purple}Cross-Validation}
\begin{itemize}
    \item Divide os dados em $k$ partes (\emph{folds})
    \item Treina em $k-1$ e testa no fold restante
    \item Repete $k$ vezes, rotacionando o fold de teste
    \item Resultado: \highlight{média ± desvio} do erro nos $k$ folds
\end{itemize}

\vspace{0.5em}
\begin{display}[Regra prática]
$k = 5$ ou $k = 10$ são os valores mais usados. Estratifique para classificação desbalanceada.
\end{display}
\end{columns}

\end{frame}


\begin{frame}[c]{$k$-Fold Cross-Validation}

\begin{figure}
    \includegraphics[height=.78\textheight]{kfold_diagram.jpg}
\end{figure}

\end{frame}


\begin{frame}[c, fragile]{Cross-Validation na prática}

\begin{code}[scikit-learn: avaliação e seleção de modelo]{python}
from sklearn.model_selection import cross_val_score, StratifiedKFold
from sklearn.ensemble import RandomForestClassifier
from sklearn.pipeline import Pipeline
from sklearn.preprocessing import StandardScaler
# Pipeline: pré-processamento + modelo em uma etapa
pipe = Pipeline([
    ("scaler", StandardScaler()),
    ("clf",    RandomForestClassifier(n_estimators=100, random_state=42)),
])
cv = StratifiedKFold(n_splits=5, shuffle=True, random_state=42)
scores = cross_val_score(pipe, X, y,
                         cv=cv, scoring="f1_weighted")
print(f"F1: {scores.mean():.3f} ± {scores.std():.3f}")
\end{code}

\begin{itemize}
    \item \texttt{Pipeline} garante que o \texttt{StandardScaler} \highlight{não veja o fold de teste}
    \item Sem Pipeline, o scaler vazaria informação do teste para o treino
\end{itemize}
\end{frame}


\begin{frame}[c, fragile]{Seleção de hiperparâmetros}

\begin{code}[GridSearchCV e RandomizedSearchCV]{python}
from sklearn.model_selection import RandomizedSearchCV
from scipy.stats import randint
param_dist = {
    "clf__n_estimators":  randint(50, 300),
    "clf__max_depth":     [None, 5, 10, 20],
    "clf__min_samples_split": randint(2, 20)}
search = RandomizedSearchCV(
    pipe,
    param_distributions=param_dist,
    n_iter=30,            # quantas combinações testar
    cv=5,
    scoring="f1_weighted",
    n_jobs=-1,            # usa todos os cores
    random_state=42)
search.fit(X_train, y_train)
print(search.best_params_)
print(f"Melhor F1 (CV): {search.best_score_:.3f}")
\end{code}

\end{frame}


\begin{frame}[c]{Variantes de Cross-Validation}

\begin{figure}
    \includegraphics[height=.78\textheight]{cv_variants.png}
\end{figure}

\end{frame}


% ============================================================
\section{HTML, CSS e a estrutura da Web}
% ============================================================

\begin{frame}[c]{Como funciona a Web}

\begin{columns}
\column{0.5\textwidth}
\textbf{\color{nes_dark_purple}O que acontece ao acessar um site}
\begin{enumerate}
    \item Navegador envia \highlight{requisição HTTP} ao servidor
    \item Servidor responde com \highlight{HTML} (estrutura)
    \item Navegador baixa \highlight{CSS} (estilo) e \highlight{JavaScript} (comportamento)
    \item Navegador \emph{renderiza} a página — monta a árvore \highlight{DOM}
\end{enumerate}

\textbf{\color{nes_dark_purple}O que é o DOM?}
\begin{itemize}
    \item \emph{Document Object Model} — representação em árvore do HTML
    \item Cada elemento HTML é um \highlight{nó} da árvore
    \item JavaScript e ferramentas de scraping navegam nessa árvore
\end{itemize}

\column{0.5\textwidth}
\textbf{\color{nes_dark_purple}HTML mínimo}
\begin{itemize}
    \item \texttt{<tag>} abre, \texttt{</tag>} fecha
    \item \texttt{id} — identificador único por página
    \item \texttt{class} — grupo de elementos com estilo comum
    \item Atributos: \texttt{href}, \texttt{src}, \texttt{data-*}
\end{itemize}

\textbf{\color{nes_dark_purple}CSS Selectors — a chave do scraping}
\begin{itemize}
    \item \texttt{div.produto} — \texttt{<div class="produto">}
    \item \texttt{\#preco} — elemento com \texttt{id="preco"}
    \item \texttt{table > tr > td} — \texttt{<td>} filho direto
\end{itemize}
\end{columns}

\end{frame}


\begin{frame}[c]{Anatomia de uma página HTML}


\begin{figure}
    \includegraphics[height=.78\textheight]{html_dom.jpg}
\end{figure}

\end{frame}


% ============================================================
\section{Web Scraping}
% ============================================================

\begin{frame}[c]{Quando o Web Scraping se torna necessário}

\begin{columns}
\column{0.5\textwidth}
\textbf{\color{nes_dark_purple}Situações típicas}
\begin{itemize}
    \item Fonte \highlight{não oferece API} pública
    \item API existe, mas tem \highlight{limites de requisição} ou paywall
    \item Dados históricos não disponíveis para download
    \item Monitoramento de preços, notícias ou vagas em tempo real
    \item Pesquisa acadêmica em dados públicos da web
\end{itemize}

\column{0.5\textwidth}
\textbf{\color{nes_dark_purple}Antes de fazer scraping, verifique}
\begin{itemize}
    \item \highlight{\texttt{robots.txt}} — o site permite crawlers?
    \item \highlight{Termos de serviço} — há proibição explícita?
    \item Existe uma \highlight{API ou dataset} oficial?
    \item Os dados são \highlight{públicos} ou requerem login?
\end{itemize}
\end{columns}

\begin{display}[Atenção legal]
Scraping de dados pessoais pode violar a LGPD. Scraping que sobrecarrega servidores pode configurar negação de serviço. Use com responsabilidade e respeite os limites.
\end{display}
\end{frame}


\begin{frame}[c, fragile]{BeautifulSoup — parsing de HTML}

\begin{code}[Instalação e uso básico]{python}
# pip install requests beautifulsoup4 lxml
import requests
from bs4 import BeautifulSoup

url = "https://pt.wikipedia.org/wiki/Inteligência_artificial"
resp = requests.get(url, headers={"User-Agent": "Mozilla/5.0"})
soup = BeautifulSoup(resp.text, "lxml")
# Título da página
print(soup.find("h1").text)

# Todos os parágrafos do artigo
paragrafos = soup.select("div.mw-parser-output > p")
for p in paragrafos[:3]:
    print(p.get_text(strip=True))
# Todos os links internos da Wikipedia
links = [a["href"] for a in soup.find_all("a", href=True)
         if a["href"].startswith("/wiki/")]
\end{code}

\end{frame}


\begin{frame}[c, fragile]{BeautifulSoup — extraindo tabelas}

\begin{code}[Tabelas HTML → DataFrame Pandas]{python}
import pandas as pd

url = "https://pt.wikipedia.org/wiki/Lista_de_pa%C3%ADses_por_PIB"
resp = requests.get(url, headers={"User-Agent": "Mozilla/5.0"})
soup = BeautifulSoup(resp.text, "lxml")

# pd.read_html usa BeautifulSoup internamente
tabelas = pd.read_html(resp.text)       # retorna lista de DataFrames
df_pib = tabelas[0]                     # primeira tabela da página
print(df_pib.head())

# Alternativa manual: navegar pela árvore
tabela = soup.find("table", class_="wikitable")
linhas = tabela.find_all("tr")
dados = [[td.get_text(strip=True) for td in tr.find_all("td")]
         for tr in linhas[1:]]          # pula cabeçalho
df = pd.DataFrame(dados, columns=["pais", "pib", "ano"])
\end{code}

\end{frame}


\begin{frame}[c]{Exemplo completo: Wikipedia}

% Sugestão de imagem: diagrama de fluxo do exemplo da Wikipedia em 4 etapas:
% (1) requests.get(url) → ícone de rede com seta para baixo
% (2) BeautifulSoup(html) → ícone de árvore DOM
% (3) soup.select / find_all → ícone de lupa
% (4) pd.DataFrame → ícone de tabela/planilha
% Cada etapa com o nome da função principal e uma linha de código abaixo.
% Layout horizontal com setas conectando as etapas.

\begin{figure}
    \includegraphics[height=.78\textheight]{scraping_flow.jpg}
\end{figure}

\end{frame}


\begin{frame}[c, fragile]{Selenium — navegadores automatizados}

Quando o conteúdo é gerado por \highlight{JavaScript} ou requer interação (login, clique, scroll), o BeautifulSoup não basta — precisamos de um navegador real.

\begin{code}[Configuração e navegação básica]{python}
# pip install selenium webdriver-manager
from selenium import webdriver
from selenium.webdriver.common.by import By
from selenium.webdriver.chrome.options import Options
opts = Options()
opts.add_argument("--headless")          # sem abrir janela
opts.add_argument("--no-sandbox")
driver = webdriver.Chrome(options=opts)
driver.get("https://quotes.toscrape.com/")
# Localizar elementos — mesmos seletores CSS
frases = driver.find_elements(By.CSS_SELECTOR, "span.text")
autores = driver.find_elements(By.CSS_SELECTOR, "small.author")
for frase, autor in zip(frases, autores):
    print(f"{autor.text}: {frase.text}")
driver.quit()
\end{code}

\end{frame}


\begin{frame}[c, fragile]{Selenium — interação e espera}

\begin{code}[Cliques, formulários e espera inteligente]{python}
from selenium.webdriver.support.ui import WebDriverWait
from selenium.webdriver.support import expected_conditions as EC
import time
driver.get("https://quotes.toscrape.com/login")
# Preencher formulário
driver.find_element(By.ID, "username").send_keys("user")
driver.find_element(By.ID, "password").send_keys("password")
driver.find_element(By.CSS_SELECTOR, "input[type='submit']").click()
# Esperar elemento aparecer (melhor que time.sleep fixo)
wait = WebDriverWait(driver, timeout=10)
wait.until(EC.presence_of_element_located(
    (By.CSS_SELECTOR, "div.quote")
))
# Scroll até o final da página (lazy loading)
driver.execute_script("window.scrollTo(0, document.body.scrollHeight);")
time.sleep(1)  # aguardar carregamento após scroll
\end{code}

\end{frame}


\begin{frame}[c]{Técnicas anti-scraping — nem tudo são flores}

\begin{columns}
\column{0.5\textwidth}
\textbf{\color{nes_dark_purple}Detecção de bots}
\begin{itemize}
    \item \highlight{Rate limiting} — bloqueia muitas requisições rápidas
    \item \highlight{CAPTCHAs} — reCAPTCHA v2/v3, hCaptcha
    \item \highlight{Fingerprinting} — tamanho de tela, plugins, fonte
    \item Detecção de \texttt{webdriver} no JavaScript
    \item Análise de padrão de comportamento (mouse, scroll)
\end{itemize}

\column{0.5\textwidth}
\textbf{\color{nes_dark_purple}Bloqueios técnicos}
\begin{itemize}
    \item \highlight{Cloudflare / WAF} — challenge automático
    \item \highlight{Conteúdo dinâmico} — dados só no JS, não no HTML
    \item \highlight{Tokens e cookies} — sessão expirada propositalmente
    \item \highlight{Honeypots} — links invisíveis que detectam bots
    \item Geração de seletores CSS aleatórios a cada deploy
\end{itemize}
\end{columns}

\begin{display}[Alternativas quando o scraping falha]
Procure APIs não documentadas (aba Network do DevTools), datasets oficiais, arquivos de exportação ou parceiros de dados. 
\end{display}

\end{frame}


% ============================================================
\section{APIs REST}
% ============================================================

\begin{frame}[c, fragile]{Consumindo APIs REST}

\begin{code}[requests — a biblioteca padrão para HTTP]{python}
import requests
# GET simples
resp = requests.get("https://api.github.com/repos/pandas-dev/pandas")
resp.raise_for_status()          # lança exceção se status >= 400
dados = resp.json()              # deserializa JSON automaticamente
print(dados["stargazers_count"], dados["forks_count"])
resp = requests.get(
    "https://api.github.com/search/repositories",
    params={"q": "language:python stars:>10000", "sort": "stars"},
    headers={"Authorization": "Bearer SEU_TOKEN"},
)
repos = resp.json()["items"]
page = 1
todos = []
while True:
    r = requests.get(url, params={"page": page, "per_page": 100})
    items = r.json()
    if not items: break
    todos.extend(items); page += 1
\end{code}

\end{frame}


\begin{frame}[c, fragile]{APIs com autenticação e rate limiting}

\begin{code}[Boas práticas: sessão, retentativas e backoff]{python}
import requests, time
from requests.adapters import HTTPAdapter
from urllib3.util.retry import Retry
# Sessão com retentativa automática
session = requests.Session()
retry = Retry(total=3,
              backoff_factor=1,          # espera 1s, 2s, 4s
              status_forcelist=[429, 500, 502, 503])
session.mount("https://", HTTPAdapter(max_retries=retry))
# Respeitar Rate Limit via header
resp = session.get("https://api.exemplo.com/dados",
                   headers={"X-API-Key": "TOKEN"})
limite  = int(resp.headers.get("X-RateLimit-Remaining", 1))
reset   = int(resp.headers.get("X-RateLimit-Reset", 0))
if limite == 0:
    espera = reset - time.time()
    print(f"Rate limit atingido. Aguardando {espera:.0f}s...")
    time.sleep(max(0, espera))
\end{code}

\end{frame}



\begin{frame}[c, noframenumbering, plain]
    \frametitle{~}
        \vfill
        \begin{center}
            {\Huge Obrigado!}\vspace{1.5em}\\
            {\Large \highlight{Alguma dúvida?}}\\
        \end{center}
        \vfill
        \begin{center}
            {\small Agora vamos para os exercícios!}
        \end{center}
\end{frame}

\end{document}